\section{Leitfaden der Usability-Tests und Abschlussinterviews}
\label{sec:BB-02_leitfaden}

\subsection{Einleitendes Moderationsskript}

\paragraph{Begrüßung}
Vielen Dank, dass Sie sich Zeit für diesen Termin nehmen.
In diesem Gespräch geht es darum, wie gut der Wiki-Auftritt Sie in Ihrem Arbeitsalltag unterstützt.
Es gibt keine richtigen oder falschen Antworten. Getestet wird nicht Ihre Leistung, sondern das System.

\paragraph{Ablauf}
Der Termin besteht aus drei Teilen:
Zunächst werden einige Hintergrundinformationen zur Rolle und zur bisherigen Nutzung des Wikis erhoben.
Anschließend bearbeiten die Teilnehmenden mehrere Aufgaben im bisherigen und im neuen Wiki und werden gebeten, ihre Gedanken laut zu äußern.
Zum Abschluss folgen Fragen zu Eindrücken und zur Akzeptanz des neuen Auftritts.

\paragraph{Hinweise zum Verhalten}
Wenn Unsicherheiten auftreten oder Teilnehmende nicht weiterwissen, ist dies besonders relevant und sollte laut geäußert werden.
Die Moderation erinnert gelegentlich an das laute Denken, greift jedoch inhaltlich möglichst wenig ein.

\paragraph{Einverständnis}
Die Ergebnisse werden ausschließlich anonymisiert ausgewertet.
Rückschlüsse auf einzelne Personen sind nicht möglich.
Nach Zustimmung beginnt die Erhebung mit kurzen Hintergrundfragen.

%-------------------------------------------------------------

\subsection{Demografie- und Hintergrundfragen (vor dem Test)}

Zur Erfassung demografischer und nutzungsbezogener Merkmale werden folgende Daten erfasst: Alter, Geschlecht, Organisationszugehörigkeit, Rolle/Funktion, Zuordnung zu Personagruppe, Nutzungshäufigkeit des bisherigen Wikis, Selbsteingeschätzte digitale Kompetenz.

%-------------------------------------------------------------

\subsection{Testaufbau}

Jede Person bearbeitet die gleichen Aufgaben einmal im alten Wiki und einmal im neuen Wiki.
Erfasst werden folgende Messgrößen: Bearbeitungszeit, Anzahl der Fehler, Anzahl der Klicks, Subjektive Schwierigkeitseinschätzung.

%-------------------------------------------------------------

\subsection{Testaufgaben}

\begin{table}[H]
  \centering
  \caption{Szenarien und Aufgaben der Usability-Tests}
  \label{tab:BB-testaufgaben}
  \begin{tabularx}{\textwidth}{c L L L}
    \toprule
    Nr.                                                                                   & Szenario & Aufgabe & Erfolgskriterium \\
    \midrule
    1                                                                                     &
    Neue Mitarbeitende möchten verstehen, wofür die Produktorganisation (PO) steht. &
    Finden Sie Informationen zum Auftrag und zur grundlegenden Struktur der PO.          &
    Korrekte Seite geöffnet                                                                                                       \\

    2                                                                                     &
    Überblick über die Rolle Release Train Engineer.                                      &
    Finden Sie die Beschreibung der Rolle Release Train Engineer.                         &
    Korrekte Seite geöffnet                                                                                                       \\

    3                                                                                     &
    Vorbereitung auf das nächste PI Planning.                                             &
    Finden Sie das Datum der nächsten PI Planning Week.                                   &
    Terminseite korrekt gefunden                                                                                                  \\

    4                                                                                     &
    Verständnis der Flow Structure in der PO.                                            &
    Finden Sie das Dokument oder die Seite zur Flow Structure.                            &
    Dokument geöffnet                                                                                                             \\

    5                                                                                     &
    Kontaktaufnahme zum LACE:HUB.                                                         &
    Finden Sie die zuständigen Ansprechpersonen für das LACE:HUB.                         &
    Seite mit Verantwortlichen geöffnet                                                                                           \\
    \bottomrule
  \end{tabularx}
\end{table}

%-------------------------------------------------------------

\subsection{Leitfaden für das Abschlussinterview}

\textbf{Block A -- Wahrgenommene Nützlichkeit und Gebrauchstauglichkeit (RQ-1)}

\begin{enumerate}
  \item \textbf{GF-1 -- Gesamtnutzen (Skala)}\\
        Inwieweit unterstützt der neue Wiki-Auftritt Ihre typischen Arbeitsaufgaben effizient?
        Skala: 1 = überhaupt nicht, 5 = in sehr hohem Maße.

  \item \textbf{HF-1 -- Begründung Gesamtnutzen}\\
        Welche konkreten Elemente tragen dazu bei, dass Aufgaben schneller oder zuverlässiger erledigt werden können?

  \item \textbf{GF-2 -- Orientierung (Skala)}\\
        Wie einfach oder schwierig war es insgesamt, sich im neuen Wiki zurechtzufinden?
        Skala: 1 = sehr schwierig, 5 = sehr einfach.

  \item \textbf{HF-2 -- \Gls{mentalemodelle} und \Gls{ia}}\\
        Gab es Situationen, in denen die Struktur Ihren Erwartungen entsprochen oder nicht entsprochen hat?
\end{enumerate}

\newpage
\textbf{Block B -- \Gls{personas}, Rollenpassung und \Gls{ux}-Methoden (RQ-2)}

\begin{enumerate}
  \setcounter{enumi}{4}
  \item \textbf{GF-3 -- Rollenpassung (Skala)}\\
        Wie gut bildet der neue Wiki-Auftritt die Informationsbedarfe Ihrer Rolle ab?
        Skala: 1 = gar nicht, 5 = sehr gut.

  \item \textbf{HF-3 -- Konkrete Strukturpassung}\\
        Wo wurde Ihre Rolle oder Personagruppe besonders berücksichtigt?
        Wo fehlte diese Passung?

  \item \textbf{OF-1 -- Wahrgenommene \Gls{ux}-Methoden}\\
        Welche Bereiche wirken speziell für Nutzergruppen entworfen, welche eher generisch?
\end{enumerate}

\textbf{Block C -- Akzeptanz, Nutzungsintention und Kommunikation (RQ-3)}

\begin{enumerate}
  \setcounter{enumi}{7}
  \item \textbf{GF-4 -- Nutzungsintention (Skala)}\\
        Wie wahrscheinlich ist eine aktive Nutzung in den nächsten drei Monaten?
        Skala: 1 = sehr unwahrscheinlich, 5 = sehr wahrscheinlich.

  \item \textbf{HF-4 -- Akzeptanztreiber und -barrieren}\\
        Welche Faktoren fördern oder hemmen die Nutzung?

  \item \textbf{HF-5 -- Kommunikations- und Einführungserwartungen}\\
        Welche Kommunikations- oder Unterstützungsformate sind notwendig?
\end{enumerate}

\textbf{Block D -- Offenes Feedback}

\begin{enumerate}
  \setcounter{enumi}{10}
  \item \textbf{OF-3 -- Offene Eindrücke}\\
        Gab es besonders positive oder negative Überraschungen?

  \item \textbf{OF-4 -- Wichtigste Verbesserung}\\
        Welche konkrete Verbesserung wünschen Sie sich für das neue Wiki?
\end{enumerate}
